\documentclass[11pt,a4paper]{article}

\usepackage{amsmath,amssymb,amsthm}
\usepackage{mathtools}
\usepackage{geometry}
\usepackage{hyperref}
\usepackage{booktabs}
\usepackage{array}
\usepackage{xcolor}
\usepackage{tikz}
\usepackage{fancyhdr}
\usepackage{titlesec}
\usepackage{tcolorbox}
\usepackage{enumitem}
\usepackage{microtype}
\usepackage{graphicx}
\usepackage{float}
\usepackage{natbib}

\geometry{margin=2.5cm, top=3cm, bottom=3cm}

\definecolor{deepblue}{RGB}{0,51,102}
\definecolor{midblue}{RGB}{0,100,180}
\definecolor{lightblue}{RGB}{220,235,250}
\definecolor{deepred}{RGB}{140,0,0}
\definecolor{boxgrey}{RGB}{245,245,248}
\definecolor{gold}{RGB}{180,140,0}

\hypersetup{
  colorlinks=true,
  linkcolor=deepblue,
  citecolor=deepred,
  urlcolor=midblue
}

\tcbuselibrary{skins,breakable}

\newtcolorbox{keybox}[1][]{
  colback=lightblue,
  colframe=deepblue,
  fonttitle=\bfseries,
  title=#1,
  breakable,
  arc=3pt
}

\newtcolorbox{defbox}[1][]{
  colback=boxgrey,
  colframe=midblue,
  fonttitle=\bfseries\color{deepblue},
  title=#1,
  breakable,
  arc=3pt
}

\newtcolorbox{resultbox}[1][]{
  colback=lightblue!40,
  colframe=gold,
  fonttitle=\bfseries\color{gold!80!black},
  title=#1,
  breakable,
  arc=3pt
}

\titleformat{\section}{\large\bfseries\color{deepblue}}{{\thesection}}{1em}{}[\color{deepblue}\titlerule]
\titleformat{\subsection}{\normalsize\bfseries\color{midblue}}{{\thesubsection}}{1em}{}
\titleformat{\subsubsection}{\normalsize\itshape\color{deepblue}}{{\thesubsubsection}}{1em}{}

\pagestyle{fancy}
\fancyhf{}
\fancyhead[L]{\small\color{deepblue}\textit{Generating Functions, Field Theory, and Analytic Combinatorics}}
\fancyhead[R]{\small\color{deepblue}\thepage}
\fancyfoot[C]{\small\color{gray}Tutorial Draft}
\renewcommand{\headrulewidth}{0.4pt}

\newtheorem{theorem}{Theorem}[section]
\newtheorem{lemma}[theorem]{Lemma}
\newtheorem{proposition}[theorem]{Proposition}
\newtheorem*{remark}{Remark}

\theoremstyle{definition}
\newtheorem{definition}[theorem]{Definition}
\newtheorem{example}[theorem]{Example}

% ---- Notation shortcuts ----
\newcommand{\GF}[1]{#1(z)}
\newcommand{\EGF}[1]{\hat{#1}(z)}
\newcommand{\R}{\mathbb{R}}
\newcommand{\C}{\mathbb{C}}
\newcommand{\Z}{\mathbb{Z}}
\newcommand{\N}{\mathbb{N}}
\newcommand{\E}{\mathbb{E}}
\newcommand{\Prob}{\mathbb{P}}
\newcommand{\ad}{a^\dagger}
\newcommand{\Hcal}{\mathcal{H}}
\newcommand{\Acal}{\mathcal{A}}
\newcommand{\Tcal}{\mathcal{T}}
\newcommand{\diff}{\,\mathrm{d}}
\newcommand{\coeff}[2]{[#1^{#2}]}
\newcommand{\abs}[1]{\left|#1\right|}
\newcommand{\set}[1]{\left\{#1\right\}}

\begin{document}

% ========== TITLE PAGE ==========
\begin{titlepage}
\centering
\vspace*{2cm}
{\Huge\bfseries\color{deepblue} Generating Functions, Field Theory,\\[0.4em] and Analytic Combinatorics}\\[1em]
{\large\color{midblue} A Unified Tutorial on the Stochastic SIR Epidemic,\\
Doi--Peliti Reaction--Diffusion Theory,\\
and the Analytic Combinatorics of Critical Exponents}\\[3cm]

\begin{tcolorbox}[colback=lightblue,colframe=deepblue,width=0.85\textwidth]
\centering\itshape
This document develops three parallel frameworks---probability generating functions
for stochastic processes, Doi--Peliti field theory for reaction--diffusion systems,
and Flajolet's analytic combinatorics---showing that all three share a common
algebraic backbone and arrive at the same critical exponents by different routes.
Every identity is motivated and every definition is given before it is used.
\end{tcolorbox}

\vfill
{\small Version 1.0 \quad|\quad March 2026}
\end{titlepage}

\tableofcontents
\clearpage

% ========== INTRODUCTION ==========
\section{Introduction and Motivation}

Three apparently distinct bodies of mathematics meet at a single generating function.

\begin{itemize}[leftmargin=2em]
  \item \textbf{Stochastic epidemiology} asks: given a random SIR epidemic,
        what is the probability of an outbreak of size $n$? The answer is encoded
        in a probability generating function (PGF) satisfying a first-order PDE
        whose \emph{characteristics are the deterministic SIR orbits}.

  \item \textbf{Doi--Peliti field theory} asks: what is the long-time density
        of particles undergoing the reaction $A+A\to\emptyset$?
        The answer emerges from a path integral whose saddle-point equation
        is identical to the mean-field rate equation.

  \item \textbf{Analytic combinatorics} asks: given a recursively defined
        class of combinatorial objects, how many exist of size $n$, and what
        is the asymptotic growth? The answer comes from the type of singularity
        of the corresponding generating function.
\end{itemize}

The central message is:

\begin{keybox}[The Unifying Principle]
A phase transition, epidemic threshold, or asymptotic exponent is a
\textbf{singularity} of a generating function. The \emph{type} of singularity
(pole, square-root branch point, logarithm) determines the universality class.
The three frameworks are three different algebraic routes to the same singularity.
\end{keybox}

\subsection{Notation and Conventions}

Throughout this document:
\begin{itemize}[leftmargin=2em]
  \item $[z^n] f(z)$ denotes the coefficient of $z^n$ in the power series $f(z)$.
  \item $f(z) \sim g(z)$ as $z\to z_*$ means $f(z)/g(z)\to 1$.
  \item $a_n \sim C \cdot \alpha^n \cdot n^\beta$ means the ratio of the two sides
        tends to 1 as $n\to\infty$.
  \item $[a,b] = ab - ba$ is the commutator of operators $a$ and $b$.
  \item $\Gamma(s)$ is the Euler gamma function, with $\Gamma(1/2)=\sqrt{\pi}$.
  \item All generating functions are \emph{formal} power series unless the context
        specifies analytic continuation.
\end{itemize}

A quick-reference table of standard results, notation, and literature is provided
in Appendix~\ref{app:table}.

\clearpage

% ========== PART I: SIR ==========
\section{Part I: The Stochastic SIR Epidemic}
\label{sec:sir}

\subsection{The Deterministic SIR System}

The Kermack--McKendrick SIR model \cite{kermack1927} partitions a population of
size $N$ into susceptible ($S$), infectious ($I$), and recovered ($R$) compartments.
The deterministic dynamics are:

\begin{equation}
  \frac{dS}{dt} = -\beta SI, \qquad
  \frac{dI}{dt} = \beta SI - \gamma I, \qquad
  \frac{dR}{dt} = \gamma I,
  \label{eq:sir_ode}
\end{equation}

where $\beta > 0$ is the infection rate and $\gamma > 0$ is the recovery rate.
The \emph{basic reproduction number} is

\begin{equation}
  R_0 = \frac{\beta N}{\gamma}.
  \label{eq:R0}
\end{equation}

\subsubsection{The Conservation Law}

Dividing the $I$-equation by the $S$-equation:

\[
  \frac{dI}{dS} = \frac{\beta SI - \gamma I}{-\beta SI} = -1 + \frac{\gamma}{\beta S}
  = -1 + \frac{1}{R_0 S/N}.
\]

Integrating (with $s = S/N$ the normalised susceptible fraction):

\begin{equation}
  \boxed{ I + S - \frac{N}{R_0}\ln S = C_0, }
  \label{eq:sir_conserved}
\end{equation}

where $C_0 = I(0) + S(0) - \frac{N}{R_0}\ln S(0)$ is fixed by initial conditions.
Equation~\eqref{eq:sir_conserved} defines the \textbf{SIR orbit} in $(S,I)$ space.
It is the single most important equation in this document; we will rederive it
from two completely different directions in Parts II and III.

\subsubsection{The Epidemic Threshold}

The epidemic peak occurs where $dI/dt = 0$, which gives $S_* = N/R_0$, i.e.\ the
threshold susceptible fraction is $s_* = 1/R_0$. The epidemic grows if $S > S_*$
(i.e.\ $R_0 > 1$) and decays if $S < S_*$.

\begin{keybox}[Epidemic Threshold]
$R_0 = 1$ is the critical point. For $R_0 > 1$ a major epidemic is possible;
for $R_0 < 1$ the epidemic always dies out. The bifurcation is continuous
(second-order in the statistical-physics sense).
\end{keybox}

\subsection{The Stochastic SIR: Master Equation}

The stochastic SIR is a continuous-time Markov chain on $(S, I)$ with transitions:

\[
  (S,I) \to (S-1,I+1) \quad\text{at rate}\quad \frac{\beta S I}{N},
  \qquad
  (S,I) \to (S, I-1) \quad\text{at rate}\quad \gamma I.
\]

Let $P(s, i, t)$ be the probability of state $(S=s, I=i)$ at time $t$.
The \textbf{master equation} (Kolmogorov forward equation) is:

\begin{align}
  \partial_t P(s,i,t) &=
    \frac{\beta(s+1)(i-1)}{N} P(s+1,i-1,t)
    + \gamma(i+1) P(s, i+1, t) \notag\\
  &\quad - \left[\frac{\beta si}{N} + \gamma i\right] P(s,i,t).
  \label{eq:master_sir}
\end{align}

This is a system of $(N+1)^2$ coupled ODEs --- impossible to solve directly
for large $N$.

\subsection{The Probability Generating Function}

\begin{definition}[Joint PGF]
The joint probability generating function of $(S, I)$ is:
\begin{equation}
  G(x, y, t) = \sum_{s=0}^{N}\sum_{i=0}^{N} P(s,i,t)\, x^s y^i.
  \label{eq:sir_pgf}
\end{equation}
\end{definition}

The dummy variables $x$ and $y$ serve purely as \emph{size markers}: the coefficient
of $x^s y^i$ in $G$ is the probability of exactly $s$ susceptibles and $i$
infectives. No probabilistic meaning is attached to $x$ or $y$ themselves.

\subsubsection{Why This Helps}

The shift operators in the master equation act as multiplication and differentiation
on $G$. For example, $\sum_{s,i} s\, P(s,i) x^s y^i = x\,\partial_x G$.
Multiplying each term in \eqref{eq:master_sir} by $x^s y^i$,
summing over all states, and using the shift identities
$\sum (s+1)P(s{+}1,i{-}1)x^s y^i = y\,\partial_x G$ and
$\sum (i+1)P(s,i{+}1)x^s y^i = \partial_y G$,
the master equation becomes a \textbf{single PDE}:

\begin{equation}
  \partial_t G = \frac{\beta}{N}(y - x)\,\partial_x(x\,\partial_y G)
                + \gamma(1 - y)\,\partial_y G.
  \label{eq:sir_pgf_pde}
\end{equation}

This is a first-order linear PDE in $(x, y, t)$ --- amenable to the method of
characteristics.

\subsection{Method of Characteristics}

\subsubsection{Setup}

For a PDE of the form $\partial_t G + A(x,y)\partial_x G + B(x,y)\partial_y G = 0$,
the \textbf{characteristics} are curves $(x(t), y(t))$ along which $G$ is constant.
They satisfy:

\[
  \frac{dx}{dt} = A(x,y), \qquad \frac{dy}{dt} = B(x,y).
\]

Reading off from \eqref{eq:sir_pgf_pde} (after rearranging to standard form):

\begin{equation}
  \frac{dx}{dt} = -\frac{\beta}{N} x(y-1), \qquad
  \frac{dy}{dt} = \frac{\beta}{N} y(y-1) - \gamma(y-1) = (y-1)\!\left(\frac{\beta y}{N} - \gamma\right).
  \label{eq:sir_chars}
\end{equation}

\subsubsection{The Key Result}

\begin{theorem}[Kendall 1956, \cite{kendall1956}]
The characteristic equations \eqref{eq:sir_chars} of the SIR PGF PDE are, under
the substitution $x = S/N$, $y = I/N$, identical to the deterministic SIR
equations \eqref{eq:sir_ode}. Therefore, the probability distribution $P(s,i,t)$
is transported along the deterministic SIR orbits in $(S,I)$ space.
\end{theorem}

\begin{proof}[Verification]
Under $x = S/N$, $\dot{x} = \dot{S}/N = -\beta SI/N^2 = -(\beta/N)x(Ni/N)$.
Comparing with \eqref{eq:sir_chars}: $\dot{x} = -(\beta/N)x(y-1)$. At leading order
in $I/N$, $y - 1 \approx y = I/N$, giving $\dot{x} = -(\beta/N)x\cdot I/N$, which
matches $\dot{S}/N$. The $y$-equation similarly recovers $\dot{I}$. \qed
\end{proof}

\begin{keybox}[The Stochastic--Deterministic Bridge]
The stochastic probability distribution $P(s,i,t)$ is constant along the
\emph{deterministic} SIR trajectories. The deterministic system is not an
approximation to the stochastic one: they share the same invariant curves exactly,
at the level of the generating function.
\end{keybox}

This result was first made explicit by Kendall \cite{kendall1956} and placed in the
PDE framework by Chalub and Souza \cite{chalub2011}. Whittle \cite{whittle1955} had
earlier used generating function methods to derive the final-size distribution.

\subsection{Critical Exponents for the SIR}

\subsubsection{Final Size Distribution}

Solving along characteristics and applying boundary conditions at $I = 0$ gives the
\textbf{final size distribution} (the probability that exactly $n$ individuals were
infected in total). For the Markovian SIR with one initial infective in a large
population:

\begin{equation}
  P(\text{total size} = n) = \frac{e^{-R_0 n}(R_0 n)^{n-1}}{n!}.
  \label{eq:sir_finalsize}
\end{equation}

This is the \textbf{Borel distribution}. It is derivable directly from Lagrange
inversion (see Part III). The generating function $T(z) = \sum_n P_n z^n$ satisfies:

\begin{equation}
  T(z) = z\, e^{R_0(T(z) - 1)}.
  \label{eq:sir_lagrange}
\end{equation}

\subsubsection{Singularity and Critical Exponent}

The singularity of $T(z)$ occurs where the implicit function theorem fails
(see Section~\ref{sec:lagrange} for the general theory). Setting
$\partial/\partial T[T - ze^{R_0(T-1)}] = 0$:

\begin{equation}
  T_* = \frac{1}{R_0}, \qquad z_* = \frac{e^{R_0 - 1}}{R_0}.
  \label{eq:sir_singularity}
\end{equation}

Near $z_*$, Taylor expanding to second order:

\begin{equation}
  T(z) \approx \frac{1}{R_0} - \sqrt{\frac{2}{R_0}} \cdot (z_* - z)^{1/2}
  + \mathcal{O}(z_* - z).
  \label{eq:sir_sqrtbranch}
\end{equation}

This is a \textbf{square-root branch point}. By the transfer theorem
(Section~\ref{sec:transfer}):

\begin{equation}
  \boxed{P(\text{size} = n) \sim C \cdot n^{-3/2} \cdot z_*^{-n}, \qquad n\to\infty.}
  \label{eq:sir_asymptotics}
\end{equation}

At criticality $R_0 = 1$: $z_* = 1$, and the $z_*^{-n}$ factor disappears, leaving
the \textbf{pure power law} $P_n \sim C n^{-3/2}$.

\begin{resultbox}[SIR Critical Exponent Summary]
\begin{center}
\begin{tabular}{lll}
\toprule
Quantity & Formula & Derivation \\
\midrule
Epidemic threshold & $R_0 = 1$ & $T_* = 1/R_0 \to 1$ \\
GF singularity & $z_* = e^{R_0-1}/R_0$ & IFT failure \\
Singularity type & Square-root branch & $F_{TT} \neq 0$ \\
Final-size decay & $P_n \sim n^{-3/2}$ at $R_0=1$ & Transfer theorem \\
Final-size GF & Borel distribution & Lagrange inversion \\
\bottomrule
\end{tabular}
\end{center}
\end{resultbox}

\clearpage

% ========== PART II: DOI-PELITI ==========
\section{Part II: Doi--Peliti Field Theory}
\label{sec:dp}

\subsection{Physical Motivation}

Consider a population of identical particles $A$ in a volume, undergoing
diffusion and the annihilation reaction $A + A \to \emptyset$ at rate $\lambda$.
In one dimension, the mean particle density $\rho(t)$ is observed to decay as:

\[
  \rho(t) \sim t^{-1/2}, \qquad t\to\infty.
\]

Deriving this rigorously --- including the universal amplitude --- requires going
beyond the mean-field rate equation $\dot\rho = -\lambda\rho^2$, which gives only
$\rho \sim (\lambda t)^{-1}$. Fluctuations dominate in low dimensions, and a
field-theoretic treatment is needed.

\subsection{Second Quantisation of the Master Equation}

\subsubsection{The Fock Space}

Let the state of the system at site $x$ be a non-negative integer $n_x$ (the
particle number). Define a \textbf{Fock space} $\mathcal{F}$ with basis vectors
$|n\rangle$ (for a single site, suppressing spatial index for now). The state of
the system at time $t$ is encoded in the \textbf{probability vector}:

\begin{equation}
  |\psi(t)\rangle = \sum_{n=0}^\infty P(n,t)\,|n\rangle.
  \label{eq:fock_state}
\end{equation}

\begin{definition}[Creation and Annihilation Operators]
Define operators $a$ (annihilation) and $\ad$ (creation) by their action on
basis states:
\begin{equation}
  a\,|n\rangle = n\,|n-1\rangle, \qquad
  \ad|n\rangle = |n+1\rangle.
  \label{eq:ladder}
\end{equation}
\end{definition}

These satisfy the canonical commutation relation:

\begin{equation}
  \boxed{[a, \ad] = 1.}
  \label{eq:ccr}
\end{equation}

This is the \textbf{Weyl algebra} (or Heisenberg algebra), identical to the
algebra of quantum mechanical ladder operators. The difference from quantum
mechanics is purely interpretational: here $|n\rangle$ is a probability state,
not a quantum state, so no Born rule applies.

The number operator $\hat{N} = \ad a$ satisfies $\hat{N}|n\rangle = n|n\rangle$.

\subsubsection{Why Annihilation and Creation?}

In the master equation, the transition $n \to n-k$ (a reaction removing $k$
particles) appears as:

\[
  \text{rate} \times P(n) \sim \text{rate} \times [n(n-1)\cdots(n-k+1)]\, P(n).
\]

The prefactor $n(n-1)\cdots(n-k+1)$ is exactly the action of $a^k$ on $|n\rangle$
(followed by appropriate factors). So the master equation becomes:

\begin{equation}
  \partial_t |\psi\rangle = \Hcal\,|\psi\rangle,
  \label{eq:doit_master}
\end{equation}

where $\Hcal$ is a polynomial in $a$ and $\ad$ called the \textbf{Liouvillian}
(or pseudo-Hamiltonian). This was Doi's key insight \cite{doi1976a, doi1976b}.

\subsubsection{The Liouvillian for $A + A \to \emptyset$}

For pair annihilation at rate $\lambda$, the transition $n \to n-2$ occurs at rate
$\lambda n(n-1)$. Writing the master equation and expressing it in terms of $a, \ad$:

\begin{equation}
  \Hcal = \lambda\bigl(1 - \ad{}^2\bigr)a^2.
  \label{eq:liouvillian_aa}
\end{equation}

\begin{remark}
The structure $(\text{final state operator} - \text{initial state operator})
\times \text{rate operator}$ is universal: any reaction $kA \to \ell A$ gives
$\Hcal = r(\ad{}^\ell - \ad{}^k)a^k$. See Appendix~\ref{app:table}.
\end{remark}

\subsection{The Generating Function as a Coherent State}

\subsubsection{Representation on Power Series}

Represent $\ad \to z$ (multiplication by $z$) and $a \to \partial_z$ (differentiation).
Then:

\begin{equation}
  |\psi(t)\rangle \equiv G(z,t) = \sum_{n=0}^\infty P(n,t)\,z^n,
  \label{eq:dp_gf}
\end{equation}

and the evolution equation \eqref{eq:doit_master} becomes:

\begin{equation}
  \partial_t G(z,t) = \Hcal(z, \partial_z)\, G(z,t).
  \label{eq:dp_gfeq}
\end{equation}

\begin{keybox}[The Fundamental Identification]
The probability generating function $G(z,t) = \sum_n P(n,t)z^n$ is \emph{identical}
to the Doi Fock-space state vector in the representation $\ad = z$, $a = \partial_z$.
The Weyl algebra acts on the space of formal power series exactly as on Fock space.
This is not an analogy: it is an isomorphism.
\end{keybox}

\subsubsection{The Weyl Algebra Explicitly}

The Weyl algebra $W_1$ is the algebra of polynomial differential operators generated
by $z$ and $\partial_z$ with relation $[\partial_z, z] = 1$. It acts on $\C[[z]]$
(formal power series) via:

\[
  z \cdot \sum_n c_n z^n = \sum_n c_n z^{n+1}, \qquad
  \partial_z \cdot \sum_n c_n z^n = \sum_n n\,c_n z^{n-1}.
\]

Every reaction Liouvillian is an element of $W_1$ (or its spatial version
$W_1 \otimes C^\infty(x)$ for field theories). The master equation is an
ODE in $W_1$.

\subsection{The Path Integral (Peliti's Construction)}

Peliti \cite{peliti1985} showed how to convert the operator evolution
\eqref{eq:doit_master} into a functional integral by inserting resolutions of
identity in coherent states at each time step. The result is:

\begin{equation}
  \langle \text{observable} \rangle =
  \int \mathcal{D}[\phi, \tilde\phi]\; \mathcal{O}[\phi]\; e^{-S[\phi,\tilde\phi]},
  \label{eq:path_integral}
\end{equation}

where the \textbf{action} for a single-species system with reaction $A+A\to\emptyset$
and diffusion is (after the \emph{Doi shift} $\tilde\phi \to 1 + \tilde\phi$):

\begin{equation}
  S = \int d^dx\,dt\; \left[
    \tilde\phi\,(\partial_t - D\nabla^2)\phi
    + 2\lambda\,\tilde\phi\phi^2
    + \lambda\,\tilde\phi^2\phi^2
  \right].
  \label{eq:action_aa}
\end{equation}

The fields $\phi(x,t)$ and $\tilde\phi(x,t)$ are the Peliti \emph{response field}
and \emph{density field} respectively. The physical density is $\rho(x,t) = \langle\phi(x,t)\rangle$.

\subsubsection{Mean-Field Limit}

Ignoring fluctuations ($\tilde\phi \to 0$, saddle-point approximation) and setting
$\phi = \rho$ uniform: only the $\tilde\phi\phi^2$ term survives in the variation,
giving $\partial_t\rho = -2\lambda\rho^2$, i.e.\ the mean-field rate equation.
The path integral saddle is the deterministic orbit, exactly as the PGF characteristic
is the SIR orbit.

\subsection{Renormalisation Group Analysis}

\subsubsection{The Ward Identity Cancellation}

For the reaction $A+A\to\emptyset$ in $d$ dimensions, the quartic vertex
$\tilde\phi^2\phi^2$ in \eqref{eq:action_aa} is exactly cancelled by a
counter-term generated by the cubic vertex $\tilde\phi\phi^2$ at one loop.
This follows from a Ward identity associated with the probability conservation
constraint $H(1, a) = 0$ (see \cite{tauber2005} for the proof). The effective action is:

\begin{equation}
  S_\mathrm{eff} = \int d^dx\,dt\; \tilde\phi\,(\partial_t - D\nabla^2)\phi,
  \label{eq:action_free}
\end{equation}

i.e.\ the theory is \emph{Gaussian} in the IR. This is an exact statement,
valid to all loop orders in $d \leq 2$.

\subsubsection{Computing the Density Decay}

From \eqref{eq:action_free}, the retarded Green's function (propagator) is:

\begin{equation}
  G_R(k, \omega) = \frac{1}{-i\omega + Dk^2}.
  \label{eq:propagator}
\end{equation}

The density at time $t$ is obtained by integrating out all wavenumber modes:

\begin{equation}
  \rho(t) = \rho_0 \int \frac{d^dk}{(2\pi)^d}\; e^{-Dk^2 t}
           = \frac{\rho_0}{(4\pi D t)^{d/2}}.
  \label{eq:density_decay}
\end{equation}

In $d = 1$:

\begin{equation}
  \boxed{\rho(t) \sim t^{-1/2}.}
  \label{eq:rho_half}
\end{equation}

The exponent $1/2$ comes from the Gaussian integral over a single momentum component
in one spatial dimension. The upper critical dimension is $d_c = 2$; for $d > d_c$
the mean-field result $\rho \sim t^{-1}$ holds.

\subsubsection{The Full Doi--Peliti Calculation: Lee (1994)}

The above treatment is schematic. The complete RG calculation for $kA\to\emptyset$
was performed by Lee \cite{lee1994}, who showed:

\begin{itemize}[leftmargin=2em]
  \item The density decay exponent $\alpha_d = d/2$ for $d < d_c = 2/(k-1)$ is
        \emph{exact} --- determined by dimensional analysis alone, not corrected
        by loops.
  \item For $k=2$ (pair annihilation), $d_c = 2$ and $\alpha_1 = 1/2$ is exact.
  \item The amplitude requires an explicit one-loop calculation.
\end{itemize}

The comprehensive review by T\"auber, Howard, and Vollmayr-Lee \cite{tauber2005}
presents the full pedagogical derivation including Feynman rules, renormalisation,
and the Callan--Symanzik equation.

\begin{resultbox}[Doi--Peliti Critical Exponent Summary]
\begin{center}
\begin{tabular}{lll}
\toprule
Quantity & Formula & Source \\
\midrule
Upper critical dimension & $d_c = 2/(k-1)$ & Dimensional analysis \\
Density decay ($d < d_c$) & $\rho \sim t^{-d/2}$ & Gaussian integral \\
Density decay ($d > d_c$) & $\rho \sim t^{-1}$ & Mean-field \\
Density decay ($d=1$, $k=2$) & $\rho \sim t^{-1/2}$ & Exact \\
Exponent exactness & All orders in $\epsilon = d_c - d$ & Ward identity \\
\bottomrule
\end{tabular}
\end{center}
\end{resultbox}

\subsection{The Linking Structure: Same Square Root}

Before proceeding to analytic combinatorics, it is worth noting the common
algebraic structure between the SIR and Doi--Peliti results.

In Laplace space ($t \to s$), the density propagator is:

\[
  \hat\rho(s) = \int \frac{d^dk}{(2\pi)^d} \frac{1}{s + Dk^2} \sim s^{-1+d/2}
  \qquad (d=1) \implies \hat\rho(s) \sim s^{-1/2}.
\]

For the SIR final-size distribution, the Laplace transform of the $n^{-3/2}$ tail is:

\[
  \sum_n n^{-3/2} e^{-sn} \sim \int_0^\infty n^{-3/2} e^{-sn}\diff n = \Gamma(-1/2)\cdot s^{1/2} \sim s^{1/2}.
\]

The survival probability $P_{\mathrm{surv}}(t) = \sum_{n>t} P_n \sim t^{-1/2}$ in Laplace space gives $\sim s^{-1/2}$.

Both are shadows of the same $\sqrt{s}$, which is the Laplace transform signature
of Brownian motion. The square root appears because the underlying conserved
quantity --- the SIR orbit \eqref{eq:sir_conserved} and the Doi diffusion
propagator \eqref{eq:propagator} --- both encode a one-dimensional Gaussian
random walk.

\clearpage

% ========== PART III: ANALYTIC COMBINATORICS ==========
\section{Part III: Analytic Combinatorics}
\label{sec:ac}

\subsection{What is Analytic Combinatorics?}

Analytic combinatorics \cite{flajolet2009} is a systematic framework for translating
between:

\begin{enumerate}[leftmargin=2em]
  \item The \textbf{recursive structure} of a combinatorial class (how its objects
        are built from smaller objects), and
  \item The \textbf{asymptotic count} of objects of size $n$.
\end{enumerate}

The bridge is the generating function. The route is:

\[
  \textbf{Grammar} \xrightarrow{\text{symbolic method}}
  \textbf{GF equation} \xrightarrow{\text{complex analysis}}
  \textbf{Asymptotics.}
\]

The power of the method is that both arrows are \emph{mechanical}: given the
grammar, the GF equation follows by a dictionary; given the GF equation, asymptotics
follow from the type of its dominant singularity.

\subsection{Combinatorial Classes and Generating Functions}

\begin{definition}[Combinatorial Class]
A \textbf{combinatorial class} is a set $\mathcal{A}$ with a \emph{size} function
$|\cdot|: \mathcal{A} \to \N$ such that the number of objects of each size is
finite. Write $a_n = |\{\alpha \in \mathcal{A} : |\alpha| = n\}|$.
\end{definition}

\begin{definition}[Ordinary Generating Function]
The \textbf{ordinary generating function (OGF)} of a class $\mathcal{A}$ is the
formal power series:
\begin{equation}
  A(z) = \sum_{n \geq 0} a_n z^n.
  \label{eq:ogf}
\end{equation}
The variable $z$ is a \emph{size marker}: the coefficient of $z^n$ counts objects
of size $n$.
\end{definition}

\begin{definition}[Exponential Generating Function]
For \emph{labelled} structures (where atoms are distinguishable), use the
\textbf{exponential generating function (EGF)}:
\begin{equation}
  \hat{A}(z) = \sum_{n \geq 0} a_n \frac{z^n}{n!}.
  \label{eq:egf}
\end{equation}
\end{definition}

\subsubsection{Why Two Kinds?}

In an unlabelled structure, atoms are indistinguishable (e.g.\ binary trees
with identical leaves). In a labelled structure, each atom carries a unique
label $1, \ldots, n$ (e.g.\ permutations, Poisson-offspring trees).

The factor $1/n!$ in the EGF precisely cancels the overcounting from relabellings.
The combinatorial constructions of Section~\ref{sec:symbolic} look identical for
both, but with different GFs.

\subsection{The Symbolic Method}
\label{sec:symbolic}

The key insight is that set-theoretic operations on combinatorial classes
correspond to algebraic operations on generating functions. The following
table holds for both OGFs (unlabelled) and EGFs (labelled):

\begin{center}
\begin{tabular}{llll}
\toprule
Construction & Notation & OGF & EGF \\
\midrule
Disjoint union & $\mathcal{A} + \mathcal{B}$ & $A(z) + B(z)$ & $\hat{A}(z) + \hat{B}(z)$ \\
Cartesian product & $\mathcal{A} \times \mathcal{B}$ & $A(z)\cdot B(z)$ & $\hat{A}(z)\cdot\hat{B}(z)$ \\
Sequence & $\mathrm{SEQ}(\mathcal{A})$ & $\frac{1}{1-A(z)}$ & $\frac{1}{1-\hat{A}(z)}$ \\
Set (unlabelled) & $\mathrm{MSET}(\mathcal{A})$ & $\exp\!\sum_{k\geq 1}\frac{A(z^k)}{k}$ & --- \\
Set (labelled) & $\mathrm{SET}(\mathcal{A})$ & --- & $e^{\hat{A}(z)}$ \\
Cycle & $\mathrm{CYC}(\mathcal{A})$ & $\log\frac{1}{1-A(z)}$ & $\log\frac{1}{1-\hat{A}(z)}$ \\
Pointing & $\Theta\mathcal{A}$ & $z\, A'(z)$ & $z\,\hat{A}'(z)$ \\
\bottomrule
\end{tabular}
\end{center}

\subsubsection{Why These Rules Hold}

\textbf{Disjoint union:} The $a_n + b_n$ objects of size $n$ in $\mathcal{A}+\mathcal{B}$
contribute $a_n z^n + b_n z^n = (a_n + b_n)z^n$.

\textbf{Cartesian product:} An object in $\mathcal{A}\times\mathcal{B}$ of size $n$
is a pair $(\alpha, \beta)$ with $|\alpha| + |\beta| = n$. Summing over all splits:
$\sum_{j=0}^n a_j b_{n-j}$, which is the coefficient of $z^n$ in $A(z)B(z)$.
This is just the Cauchy product of series.

\textbf{Sequence:} $\mathrm{SEQ}(\mathcal{A}) = \epsilon + \mathcal{A} + \mathcal{A}^2 + \ldots$
where $\epsilon$ is the empty sequence. Translating: $1 + A + A^2 + \ldots = 1/(1-A)$.

\textbf{Labelled set:} A set of $k$ labelled objects from $\mathcal{A}$ has EGF
$\hat{A}(z)^k/k!$ (divide by $k!$ because the set has no order). Summing over $k$:
$\sum_{k\geq 0} \hat{A}^k/k! = e^{\hat{A}(z)}$.

\begin{example}[Binary Trees]
A binary tree $\mathcal{T}$ is either a leaf $\bullet$ (size 1) or a pair of
binary trees joined at a root (size $= 1 + $ size of left subtree $+$ size of
right subtree). The grammar is:
\[
  \mathcal{T} = \{\bullet\} + \bullet \times \mathcal{T} \times \mathcal{T}.
\]
Translating: $T(z) = z + z\cdot T(z)^2$. Solving: $T(z) = (1-\sqrt{1-4z})/2$.
Expanding: $[z^n]T(z) = C_{n-1}$ where $C_n = \binom{2n}{n}/(n+1)$ is the $n$-th
Catalan number.
\end{example}

\subsection{Singularity Analysis}
\label{sec:singularity}

\subsubsection{Why Singularities Determine Asymptotics}

The coefficients of a GF are recovered by Cauchy's integral formula:

\begin{equation}
  a_n = [z^n]\,A(z) = \frac{1}{2\pi i}\oint_{|z|=r} \frac{A(z)}{z^{n+1}}\diff z.
  \label{eq:cauchy}
\end{equation}

For large $n$, this contour integral is dominated by the \textbf{singularities of
$A(z)$ closest to the origin} (smallest $|z_*|$). All the large-$n$ information
is compressed into the nature of these singularities.

\subsubsection{The Transfer Theorems}

Let $\rho = 1/|z_*|$ be the exponential growth rate and write $A(z) \sim K(1-z/z_*)^\alpha$
near $z_*$.

\begin{theorem}[Transfer Theorem, Flajolet--Odlyzko \cite{flajolet2009}]
If $A(z)$ is analytic in a \emph{dented domain} around $z_*$ (see Fig.~below)
and $A(z) \sim K(1 - z/z_*)^\alpha$ as $z\to z_*$ with $\alpha \notin \N$, then:

\begin{equation}
  [z^n]\,A(z) \sim \frac{K}{\Gamma(-\alpha)}\cdot n^{-\alpha-1} \cdot z_*^{-n}.
  \label{eq:transfer}
\end{equation}
\end{theorem}

\noindent The key cases:

\begin{center}
\begin{tabular}{lll}
\toprule
Singularity type & Local form & $a_n$ asymptotics \\
\midrule
Simple pole & $(1 - z/z_*)^{-1}$ & $C \cdot z_*^{-n}$ \\
$k$-th order pole & $(1 - z/z_*)^{-k}$ & $C \cdot n^{k-1} \cdot z_*^{-n}$ \\
Square-root branch & $(1-z/z_*)^{1/2}$ & $C \cdot n^{-3/2} \cdot z_*^{-n}$ \\
Inverse square-root & $(1-z/z_*)^{-1/2}$ & $C \cdot n^{-1/2} \cdot z_*^{-n}$ \\
Logarithm & $\log(1-z/z_*)$ & $C \cdot n^{-1} \cdot z_*^{-n}$ \\
General $\alpha$ & $(1-z/z_*)^\alpha$ & $C \cdot n^{-\alpha-1} \cdot z_*^{-n}$ \\
\bottomrule
\end{tabular}
\end{center}

The $\Gamma(-\alpha)$ in \eqref{eq:transfer} is the normalisation from the
generalised binomial expansion: $(1-u)^\alpha = \sum_n \binom{\alpha}{n}(-u)^n$
and $\binom{\alpha}{n} \sim n^{-\alpha-1}/\Gamma(-\alpha)$.

\subsubsection{The Dented Domain (Briefly)}

The transfer theorem requires that $A(z)$ can be analytically continued to a
region $\Delta$ that \emph{dents into} the disk $|z| < |z_*|$ from outside
the singularity. This excludes cases where there are infinitely many singularities
on the circle $|z| = |z_*|$ (which can happen for combinatorial GFs with
periodic structure). For all the GFs in this document, the condition is satisfied.

\subsection{Lagrange Inversion}
\label{sec:lagrange}

\subsubsection{The Problem}

Many combinatorial recursions take the \textbf{Lagrange form}:

\begin{equation}
  T(z) = z\,\phi(T(z)),
  \label{eq:lagrange_eq}
\end{equation}

for some known function $\phi$. The GF $T(z)$ is implicitly defined and generally
has no closed form. Lagrange (1770) gave a direct formula for its coefficients.

\subsubsection{Origin: Fixed-Point Inversion}

Equation \eqref{eq:lagrange_eq} says that $T$ is a fixed point of the map
$T \mapsto z\phi(T)$. For small $z$, this map is a contraction and the fixed point
is unique and analytic. As $z$ increases, the map eventually ceases to be a
contraction --- this is exactly the singularity.

More precisely, the implicit function theorem guarantees that $T(z)$ is analytic
near $z = 0$ as long as $\partial/\partial T[T - z\phi(T)] \neq 0$, i.e.\ as long as
$1 \neq z\phi'(T)$.

\begin{theorem}[Lagrange Inversion Formula]
If $T = z\phi(T)$ and $\phi(0) \neq 0$, then:

\begin{equation}
  [z^n]\,T(z) = \frac{1}{n}\,[T^{n-1}]\,\phi(T)^n.
  \label{eq:lagrange_inversion}
\end{equation}
\end{theorem}

\begin{proof}[Proof sketch]
Write $z = T/\phi(T)$ so that $dz = (1/\phi - T\phi'/\phi^2)dT$. By Cauchy's formula:
\[
[z^n]T = \frac{1}{2\pi i}\oint \frac{T}{z^{n+1}}\diff z
= \frac{1}{2\pi i}\oint \frac{T}{\left(T/\phi(T)\right)^{n+1}}\frac{d z}{d T}\diff T.
\]
Computing $dz/dT$, using residue calculus at $T=0$, and simplifying yields
\eqref{eq:lagrange_inversion}. \qed
\end{proof}

\subsubsection{The Singularity of the Lagrange Equation}

Setting $\partial F/\partial T = 0$ where $F(T,z) = T - z\phi(T)$:

\begin{equation}
  1 = z_*\phi'(T_*), \qquad T_* = z_*\phi(T_*).
  \label{eq:lagrange_singular}
\end{equation}

These two equations determine the branch point $(T_*, z_*)$. Near $(T_*, z_*)$,
expanding to second order in $\delta T = T - T_*$:

\[
  0 = F_z(z_* - z) + \tfrac{1}{2}F_{TT}\,(\delta T)^2 + \ldots
\]

Since $F_z = -\phi(T_*)$ and $F_{TT} = -z_*\phi''(T_*)$:

\begin{equation}
  \delta T \sim \sqrt{\frac{2\phi(T_*)}{z_*\phi''(T_*)}}\cdot(z_* - z)^{1/2}.
  \label{eq:lagrange_sqrtbranch}
\end{equation}

\textbf{Any Lagrange equation with $\phi'' \neq 0$ at the critical point produces
a square-root branch point.} The exponent $-3/2$ in the coefficient asymptotics
(via the transfer theorem) is therefore \emph{universal} for this class of equations.

\clearpage

% ========== PART IV: AC RECOVERS SIR AND DOI-PELITI ==========
\section{Part IV: Analytic Combinatorics Recovers Both Results}
\label{sec:acrecovery}

\subsection{AC Derivation of the SIR Final-Size Distribution}

\subsubsection{The Epidemic as a Labelled Tree}

In the large-$N$ limit, each infectious individual independently infects a
Poisson$(R_0)$ number of others before recovering. Because each new infection
is of a (hitherto) susceptible individual chosen uniformly at random from a
large pool, infections are approximately independent --- this is the
\emph{branching process approximation}.

The epidemic is therefore a \textbf{Galton--Watson tree} with Poisson$(R_0)$
offspring. The total number of infected individuals (epidemic size) is the
\textbf{total progeny} of this tree.

The symbolic description is: an epidemic tree $\mathcal{T}$ consists of a
labelled root followed by a labelled set of sub-epidemics (one per infection):

\begin{equation}
  \mathcal{T} = \bullet \times \mathrm{SET}(\mathcal{T}) \quad\text{(weighted by Poisson)}.
  \label{eq:sir_grammar}
\end{equation}

In EGF language, with the Poisson weight $e^{R_0(T-1)}$:

\begin{equation}
  T(z) = z\,e^{R_0(T(z) - 1)}.
  \label{eq:sir_lagrange2}
\end{equation}

This is the Lagrange equation \eqref{eq:lagrange_eq} with $\phi(T) = e^{R_0(T-1)}$.

\subsubsection{Applying Lagrange Inversion}

With $\phi(T) = e^{R_0(T-1)}$:

\[
  [z^n]T = \frac{1}{n}[T^{n-1}]\,e^{R_0 n(T-1)}
          = \frac{e^{-R_0 n}}{n}\cdot\frac{(R_0 n)^{n-1}}{(n-1)!}
          = \frac{e^{-R_0 n}(R_0 n)^{n-1}}{n!}.
\]

This recovers the Borel distribution \eqref{eq:sir_finalsize} in one line.

\subsubsection{Finding the Singularity via the Lagrange Conditions}

Applying \eqref{eq:lagrange_singular} with $\phi(T) = e^{R_0(T-1)}$,
$\phi'(T) = R_0 e^{R_0(T-1)}$:

\begin{align}
  1 &= z_* R_0 e^{R_0(T_*-1)} \label{eq:sir_crit1}\\
  T_* &= z_* e^{R_0(T_*-1)}. \label{eq:sir_crit2}
\end{align}

Dividing: $T_* R_0 = 1$, giving $T_* = 1/R_0$. Substituting back:
$z_* = (1/R_0)/e^{R_0(1/R_0 - 1)} = (1/R_0)/e^{1-R_0} = e^{R_0-1}/R_0$, confirming
\eqref{eq:sir_singularity}.

\subsubsection{The Singular Locus Is the SIR Orbit}

At the critical point, take logarithms of \eqref{eq:sir_crit2}:

\[
  \ln T_* = \ln z_* + R_0(T_* - 1).
\]

Rearranging:

\begin{equation}
  \ln z_* = \ln T_* - R_0 T_* + R_0.
  \label{eq:singular_locus}
\end{equation}

Now let $T_*$ vary continuously (treating the singularity as a function of
the critical point). Setting $T = S/N$ (susceptible fraction):

\begin{equation}
  \ln z = \ln\frac{S}{N} - R_0\frac{S}{N} + \text{const}.
  \label{eq:singular_locus_S}
\end{equation}

Compare with the SIR conservation law \eqref{eq:sir_conserved}: $I + S - (N/R_0)\ln S = C_0$,
which rearranges to:

\begin{equation}
  \ln S - R_0\frac{S}{N} = \tilde{C}.
\end{equation}

\textbf{These are identical.} The singular locus of the GF $T(z)$ traces the
SIR orbit in $(S/N, z)$ space.

\begin{keybox}[AC Recovers the SIR Conserved Quantity]
The SIR conservation law $I + S - (N/R_0)\ln S = C$ is not an input to the
analytic combinatorics calculation. It emerges as the \textbf{singular locus
of the epidemic-tree generating function} under the substitution $T = S/N$.
The method of characteristics (Part I) and the Lagrange singularity condition
(Part III) are two different derivations of the same algebraic curve.
\end{keybox}

\subsubsection{The Two Branches = Minor vs Major Epidemic}

Near $z_*$, \eqref{eq:lagrange_sqrtbranch} gives two branches $T = T_* \pm C\sqrt{z_* - z}$:

\begin{itemize}[leftmargin=2em]
  \item \textbf{Small branch} ($T < T_* = 1/R_0$): the epidemic is small, dies out
        before crossing the peak susceptible threshold. This is the \emph{minor outbreak}.
  \item \textbf{Large branch} ($T > T_* = 1/R_0$): the epidemic crosses the peak
        and infects $O(N)$ individuals. This is the \emph{major epidemic}.
\end{itemize}

The extinction probability $q$ (probability of minor outbreak) is the value of
$T(1)$ on the small branch: it satisfies $q = e^{R_0(q-1)}$, the classic
Galton--Watson extinction equation.

\subsection{AC Derivation of the $A+A\to\emptyset$ Exponent}

\subsubsection{Reduction to First Passage}

In one spatial dimension, particles cannot cross each other. Label them $1, 2, \ldots$
from left to right. Particle $k$ is annihilated when it meets particle $k+1$ (or $k-1$).
In \emph{relative} coordinates, the pair $(k, k+1)$ is described by a single
random walk $X_t = x_{k+1}(t) - x_k(t)$ on the positive integers, starting at
some $d > 0$. The pair annihilates when $X_t$ first hits 0.

The density $\rho(t)$ decays because pairs annihilate; the annihilation rate is
the first-passage rate of the relative random walk to 0. Specifically:

\begin{equation}
  \rho(t) \sim P(\text{relative walk has not hit 0 by time } t)
          = P_{\mathrm{surv}}(t).
  \label{eq:rho_surv}
\end{equation}

\subsubsection{The First-Passage GF}

A \textbf{first-passage path of length $n$} from $d > 0$ to 0 is a nearest-neighbour
path that stays strictly positive for steps $1, \ldots, n-1$ and reaches 0 at step $n$.

For the \emph{symmetric} random walk (equal left/right probabilities $p = 1/2$),
the GF for first-passage paths starting at $d = 1$ satisfies the functional equation:

\begin{equation}
  F(z) = \frac{z}{2} + \frac{z}{2}\cdot F(z)^2,
  \label{eq:firstpassage_eq}
\end{equation}

since each path either takes one step directly to 0 (weight $z/2$), or takes one
step right, executes a complete excursion back to 1, and then completes a first
passage from 1 to 0 (weight $(z/2)\cdot F(z)^2$).

\begin{remark}
Equation \eqref{eq:firstpassage_eq} is a Lagrange equation! With $p=1/2$ and
$F = z\phi(F)$ it reads $F(z) = z \cdot \frac{1+F^2}{2}$, so $\phi(F) = (1+F^2)/2$.
\end{remark}

Solving the quadratic:

\begin{equation}
  F(z) = \frac{1 - \sqrt{1 - z^2}}{z}.
  \label{eq:firstpassage_soln}
\end{equation}

\subsubsection{The Singularity}

$F(z)$ has a square-root branch point at $z_* = 1$ (where $1 - z^2 = 0$). Near $z = 1$:

\begin{equation}
  F(z) \approx 1 - \sqrt{2(1-z)} + \mathcal{O}(1-z),
  \label{eq:fp_branch}
\end{equation}

i.e.\ $F(z) \sim 1 - \sqrt{2}\cdot(1-z)^{1/2}$. With $\alpha = 1/2$, the transfer
theorem \eqref{eq:transfer} gives:

\begin{equation}
  f_n = [z^n]F(z) \sim \frac{\sqrt{2}}{2}\cdot\frac{1}{\Gamma(-1/2)}\cdot n^{-3/2}
      = \frac{1}{\sqrt{2\pi}}\cdot n^{-3/2}.
  \label{eq:fp_asymp}
\end{equation}

(Using $\Gamma(-1/2) = -2\sqrt{\pi}$.)

\subsubsection{Integrating to the Survival Probability}

The survival probability is the complementary CDF of the first-passage distribution:

\begin{equation}
  P_{\mathrm{surv}}(t) = \sum_{n > t} f_n \sim \int_t^\infty \frac{dn}{\sqrt{2\pi}\, n^{3/2}}
  = \sqrt{\frac{2}{\pi t}}.
  \label{eq:surv_asymptotics}
\end{equation}

Therefore, from \eqref{eq:rho_surv}:

\begin{equation}
  \boxed{\rho(t) \sim \left(\frac{2}{\pi t}\right)^{1/2} \sim t^{-1/2},}
  \label{eq:rho_ac}
\end{equation}

in exact agreement with the Doi--Peliti result \eqref{eq:rho_half}.

\begin{keybox}[AC Recovers the Doi--Peliti Exponent]
The exponent $-1/2$ comes from integrating the first-passage tail $f_n \sim n^{-3/2}$.
The $n^{-3/2}$ comes from the transfer theorem applied to the square-root branch
point of $F(z)$ at $z_* = 1$. The square-root branch point is forced by the
Lagrange structure of the first-passage recursion. The chain is:
\[
  \text{Annihilation} \to \text{First-passage problem} \to
  \text{Lagrange GF} \to \text{Square-root branch} \to
  n^{-3/2} \to \rho(t) \sim t^{-1/2}.
\]
\end{keybox}

\subsection{The Common Algebraic Root}

Both derivations hinge on a \textbf{square-root branch point} in a GF satisfying a
Lagrange equation. This is not a coincidence. The Lagrange equation encodes
\emph{branching structure}; and branching structure appears in both settings because:

\begin{itemize}[leftmargin=2em]
  \item In the SIR, the epidemic is a Galton--Watson branching tree.
  \item In the annihilation problem, the first-passage path is built from
        excursions which are themselves branching trees (Motzkin paths / ballot paths).
\end{itemize}

In both cases, the \emph{critical point} is where the average branching number
equals one: $R_0 = 1$ for SIR, and the random walk at criticality (equal left/right
probabilities) for first passage. At this criticality, the Lagrange singularity
sits on the boundary of the unit disk ($z_* = 1$), and the asymptotics are pure
power laws without exponential factors.

The Doi--Peliti calculation arrives at the same $t^{-1/2}$ via a completely
different path: Gaussian integration over diffusion modes in momentum space.
The Laplace-space correspondence

\[
  \int \frac{d^1k}{2\pi} \frac{1}{s + Dk^2} \sim s^{-1/2}
  \quad\longleftrightarrow\quad
  \sum_n n^{-3/2} e^{-sn} \sim s^{1/2}
\]

identifies the two computations as Fourier/Laplace duals. The combinatorial
$n^{-3/2}$ and the field-theoretic $k^{-1}$ momentum integral are the same
square root in different basis representations.

\clearpage

% ========== PART V: AC AND FEYNMAN DIAGRAMS ==========
\section{Part V: Analytic Combinatorics and Feynman Diagrams}
\label{sec:feynman}

\begin{keybox}[The Thesis of This Section]
The connection between AC and Feynman diagrams is not analogical.
It is a literal identification of the same algebraic structures
appearing simultaneously in combinatorics and quantum field theory.
There are four distinct levels at which this identification holds,
ranging from classical folklore to active research.
We develop each in turn, with explicit examples, and close with
an assessment of what is established and what remains open.
\end{keybox}

\subsection{Connection 1: The Exponential Formula and $\log Z$}
\label{sec:expformula}

\subsubsection{The Physics Statement}

In any quantum field theory or statistical mechanics model, the partition
function $Z$ receives contributions from all Feynman diagrams (connected or not).
The \textbf{free energy} $F = -\log Z$ receives contributions only from
\emph{connected} diagrams. This is used in every QFT textbook to discard
disconnected vacuum bubbles.

\subsubsection{The AC Statement}

In analytic combinatorics, if $\mathcal{C}$ is a labelled class of
\emph{connected} structures, and $\mathcal{A} = \mathrm{SET}(\mathcal{C})$
is the class of all structures (each being a set of connected components), then:

\begin{equation}
  \hat{A}(z) = e^{\hat{C}(z)}, \qquad \hat{C}(z) = \log \hat{A}(z).
  \label{eq:expformula}
\end{equation}

This is the \textbf{exponential formula}. It follows immediately from the
EGF rule for the SET construction: $\mathrm{SET}(\mathcal{C})$ translates
to $e^{\hat{C}(z)}$.

\subsubsection{They Are the Same Identity}

The statement ``$\log Z$ generates connected diagrams'' is
\emph{exactly} \eqref{eq:expformula} with:
\begin{align*}
  \hat{A}(z) &\leftrightarrow Z \quad \text{(all diagrams, vacuum normalised)}, \\
  \hat{C}(z) &\leftrightarrow \log Z \quad \text{(connected diagrams only)}, \\
  z^n/n!     &\leftrightarrow g^n \quad \text{(perturbative expansion in coupling $g$)}.
\end{align*}

The \emph{symmetry factors} of Feynman diagrams --- those awkward denominators
$1/|\mathrm{Aut}(\Gamma)|$ that appear in loop calculations --- are exactly the
$1/n!$ overcounting factors in the EGF SET construction. They have a purely
combinatorial origin: the automorphism group of a diagram is the symmetry group
that acts on its legs and propagators, which is exactly the symmetry group
counted by the EGF.

\begin{example}[Symmetry Factor from the SET Construction]
A diagram $\Gamma$ consisting of two identical one-loop bubbles (a ``figure-eight''
vacuum graph) has $|\mathrm{Aut}(\Gamma)| = 8$ in $\phi^4$ theory. In EGF language:
\[
  \hat{C}_{\text{two-bubble}}(z) = \frac{1}{2!}\hat{C}_{\text{one-bubble}}(z)^2
  = \frac{1}{2}\hat{C}_{\text{one-bubble}}(z)^2.
\]
The $1/2!$ is the SET factor for a set of two identical objects. More generally,
for a vacuum graph consisting of $k$ copies of the same connected diagram $\gamma$:
\[
  \text{symmetry factor} = |\mathrm{Aut}(\gamma)|^k \cdot k!,
\]
where $k!$ is from the EGF SET construction and $|\mathrm{Aut}(\gamma)|^k$
accounts for the internal symmetries of each copy.
\end{example}

\begin{resultbox}[Status]
Completely classical. Appears in Flajolet \& Sedgewick \cite{flajolet2009}
as a standard EGF application, and in every QFT textbook (e.g.\
Zinn-Justin \cite{zinjustin2002}) as the cluster expansion / cumulant theorem.
The explicit identification of symmetry factors with EGF overcounting is
noted in Cvitanovi\'c, Lautrup \& Pearson (1978) and in Yeats \cite{yeats2017}.
\end{resultbox}

\subsection{Connection 2: Dyson--Schwinger Equations Are Lagrange Equations}
\label{sec:dse_lagrange}

\subsubsection{The Dyson--Schwinger Equation}

In $\phi^3$ theory, the full (dressed) propagator $G(p^2)$ satisfies the
Dyson--Schwinger equation:

\begin{equation}
  G = G_0 + G_0 \cdot \Sigma(G) \cdot G,
  \label{eq:dse}
\end{equation}

where $G_0 = 1/(p^2 + m^2)$ is the bare propagator and $\Sigma(G)$ is the
self-energy, which is itself a function of $G$ (all 1PI insertions).
Rearranging:

\begin{equation}
  G = G_0\bigl[1 + \Sigma(G) \cdot G\bigr] \equiv G_0 \cdot \Phi(G).
  \label{eq:dse_lagrange}
\end{equation}

This is exactly a \textbf{Lagrange equation} $T = z\phi(T)$ with:
\[
  T \leftrightarrow G, \quad z \leftrightarrow G_0, \quad \phi(T) = 1 + \Sigma(T)\cdot T.
\]

\subsubsection{Lagrange Inversion Generates the Perturbative Expansion}

Applying \eqref{eq:lagrange_inversion}:
\begin{equation}
  [G_0^n]\, G = \frac{1}{n}\,[G^{n-1}]\,\Phi(G)^n.
  \label{eq:dse_li}
\end{equation}

Each term $[G^{n-1}]\Phi(G)^n$ counts the number of Feynman diagrams with
$n$ propagators, weighted by their amplitudes. The Lagrange inversion formula
is therefore a \textbf{generating function for the perturbative expansion} of
the dressed propagator.

\subsubsection{The Singularity is the Landau Pole}

By the Lagrange singularity conditions \eqref{eq:lagrange_singular}:
\[
  1 = G_0^* \Phi'(G^*), \qquad G^* = G_0^* \Phi(G^*).
\]

The singularity of $G$ (as a function of $G_0$ or equivalently of the coupling
$\lambda$) is the \textbf{Landau pole}: the scale at which the perturbative
expansion breaks down and non-perturbative effects take over. In AC terms, it
is the branch point of the Lagrange GF. The type of singularity (square root,
pole, essential) determines the non-perturbative physics.

\begin{example}[Running Coupling in QED]
In QED, the running coupling $\alpha(\mu^2)$ satisfies a Dyson--Schwinger-type
equation that, in the leading-log approximation, reduces to:
\[
  \alpha(\mu^2) = \frac{\alpha_0}{1 - (\alpha_0/3\pi)\log(\mu^2/m^2)}.
\]
This has a simple pole at $\mu^2 = m^2 e^{3\pi/\alpha_0}$ --- the Landau pole.
In Lagrange language, $\phi(T) = 1/(1-cT)$ produces a simple pole at
$T_* = 1/c$, giving $z_* = 1/(ce)$ and $G \sim (G_0^* - G_0)^{-1}$ near
the singularity. The transfer theorem then gives $[G_0^n]G \sim C \cdot (G_0^*)^{-n}$
--- pure exponential growth in loop order, with no power-law correction.
\end{example}

\begin{resultbox}[Status]
The identification of Dyson--Schwinger equations with fixed-point / Lagrange
equations is known in the Connes--Kreimer community
(Foissy \cite{foissy2007}, Yeats \cite{yeats2017}).
The explicit phrasing in terms of the AC Lagrange inversion theorem
with the IFT singularity condition and transfer theorem is not, as far as
we are aware, stated cleanly in the literature. The two communities (AC and
QFT) use different language for the same object.
\end{resultbox}

\subsection{Connection 3: The Connes--Kreimer Hopf Algebra}
\label{sec:ck}

This is the deepest and most surprising connection. It was a genuine
mathematical discovery, not an observation, and it resolved a long-standing
conceptual puzzle about the structure of renormalisation.

\subsubsection{The Problem: Nested Subdivergences}

In QFT, a Feynman diagram $\Gamma$ may contain divergent subdiagrams
$\gamma \subset \Gamma$. To renormalise $\Gamma$, one must subtract the
subdivergences \emph{before} subtracting the overall divergence. For
deeply nested diagrams this leads to the complicated BPHZ forest formula
(Bogoliubov--Parasiuk--Hepp--Zimmermann):

\begin{equation}
  R(\Gamma) = \Gamma + \sum_{\text{forests } \mathcal{F}} \prod_{\gamma \in \mathcal{F}}
  (-C(\gamma)) \cdot \Gamma/\mathcal{F},
  \label{eq:bphz}
\end{equation}

where $C(\gamma)$ is the counterterm for subdiagram $\gamma$ and
$\Gamma/\mathcal{F}$ is the quotient graph. The sum is over all admissible
forests of subdivergent subgraphs.

\subsubsection{The Discovery: This Is a Hopf Algebra Antipode}

Connes and Kreimer \cite{conneskreimer1998} showed that \eqref{eq:bphz}
is exactly the \textbf{antipode} $S$ of a Hopf algebra $H_R$ generated
by \emph{rooted trees}, where:

\begin{itemize}[leftmargin=2em]
  \item Each 1PI Feynman diagram $\Gamma$ maps to a rooted tree $t_\Gamma$
        whose nodes are the divergent subgraphs and whose root is $\Gamma$ itself.
  \item The \textbf{coproduct} $\Delta(t) = \sum_c P_c(t) \otimes R_c(t)$
        (summing over admissible cuts $c$) decomposes the tree into a
        subtree $P_c(t)$ (the subdivergence) and the quotient $R_c(t)$
        (the remaining diagram after subtraction).
  \item The \textbf{antipode} $S(t) = -t - \sum_{c \neq \emptyset} S(P_c(t))\cdot R_c(t)$
        is the BPHZ counterterm.
\end{itemize}

\subsubsection{The AC Identification}

In AC language:
\begin{itemize}[leftmargin=2em]
  \item The coproduct is an \textbf{admissible cut} of a rooted tree, which
        is exactly the decomposition used in the Lagrange analysis of tree GFs
        (the subtree above a cut corresponds to $P_c$; the base tree below
        corresponds to $R_c$).
  \item The antipode formula is the \textbf{M\"obius inversion} / inclusion--exclusion
        that appears in the MSET and CYC constructions in AC.
  \item The \textbf{beta function} (the RG flow equation $\mu\partial_\mu g = \beta(g)$)
        is a \textbf{derivation} of the Hopf algebra $H_R$, and its perturbative
        expansion is a GF of rooted trees in the AC sense.
\end{itemize}

\subsubsection{An Explicit Small Example}

Consider the simplest non-trivial case: a one-loop self-energy subgraph $\gamma$
inserted into a two-loop diagram $\Gamma$. The rooted tree is:

\[
  t_\Gamma = \begin{array}{c}
    \Gamma \\
    | \\
    \gamma
  \end{array}
\]

The coproduct (the only non-trivial admissible cut is at the edge $\Gamma\text{---}\gamma$):
\[
  \Delta(t_\Gamma) = t_\Gamma \otimes 1 + 1 \otimes t_\Gamma + \gamma \otimes (\Gamma/\gamma).
\]

The antipode:
\[
  S(t_\Gamma) = -t_\Gamma - S(\gamma)\cdot(\Gamma/\gamma) = -t_\Gamma + \gamma\cdot(\Gamma/\gamma).
\]

This is exactly the BPHZ formula: subtract the subdivergence $\gamma$ from
$\Gamma/\gamma$, then subtract the result from $\Gamma$.

In AC terms: the GF for the class of rooted trees of depth 2 is
$T_2(z) = z \cdot T_1(z)$ where $T_1(z) = z/(1-z)$, and the admissible
cut is a coefficient extraction in the composition.

\subsubsection{The Riemann--Hilbert Connection}

Connes and Kreimer further showed \cite{conneskreimer2000} that the
Birkhoff decomposition of a loop $\gamma: S^1 \to G$ in the group $G$
of characters of $H_R$ is the mathematical statement of renormalisation:
the ``positive part'' of the decomposition is the renormalised amplitude
and the ``negative part'' is the counterterm. This connects renormalisation
to the \textbf{Riemann--Hilbert problem} --- a deep link between physics
and classical complex analysis.

\begin{resultbox}[Status]
Rigorously established. Original papers: Connes \& Kreimer
\cite{conneskreimer1998, conneskreimer2000}. The combinatorial structure
of the Hopf algebra (without the physics) was known earlier from
Grossman--Larson (1989) in the context of numerical analysis (Runge--Kutta
methods). The book by Yeats \cite{yeats2017} is the most accessible
treatment oriented toward combinatorialists. Manchon \cite{manchon2004}
provides a concise mathematical survey.
\end{resultbox}

\subsection{Connection 4: Factorial Divergence, Borel Summation,\\
and the Transfer Theorem}
\label{sec:borel}

\subsubsection{The Problem: Perturbation Theory Always Diverges}

The number of Feynman diagrams at loop order $n$ grows as $n!$ for
large $n$. This is a combinatorial fact: the number of ways to connect
$2n$ legs in pairs (Wick contractions) is $(2n-1)!! \sim \sqrt{2}(2n/e)^n$,
which grows faster than any exponential. Therefore the perturbative expansion:
\[
  Z(g) = \sum_{n=0}^\infty c_n g^n, \qquad c_n \sim K \cdot n! \cdot \alpha^{-n},
\]
has \emph{zero radius of convergence} for any $\alpha > 0$. This is not a
defect of QFT; it is a fundamental feature.

\subsubsection{The AC Explanation}

The transfer theorem \eqref{eq:transfer} immediately explains this. If $Z(g)$
has a singularity at $g_* = 1/\alpha$ of the form:

\[
  Z(g) \sim K(1 - g/g_*)^{-\beta}, \qquad \beta \in \N,
\]

then:
\[
  c_n = [g^n]Z(g) \sim \frac{K}{\Gamma(\beta)}\cdot n^{\beta-1} \cdot g_*^{-n}.
\]

For a simple pole ($\beta = 1$): $c_n \sim K \cdot g_*^{-n}$, exponential growth.
For the partition function of $\phi^4$ theory, the singularity at finite $g_*$
is associated with a \textbf{instanton} (a non-perturbative saddle point), and:
\[
  c_n \sim K \cdot n! \cdot g_*^{-n}.
\]

The $n!$ growth corresponds to an \textbf{essential singularity} in the
Borel plane, not a branch point or pole in $g$-space. In the Borel transform
$\hat{Z}(t) = \sum_n (c_n/n!) t^n$, the $n!$ growth means the series converges
(radius of convergence $|t| < g_*$), and $\hat{Z}(t)$ has a singularity at
$t = g_*$. The Borel sum $Z(g) = \int_0^\infty e^{-t/g}\hat{Z}(t)\diff t$
is obstructed by this singularity.

\begin{keybox}[The Transfer Theorem in the Borel Plane]
The type of singularity of $\hat{Z}(t)$ at $t = g_*$ determines the
non-perturbative physics:
\begin{itemize}[leftmargin=1.5em]
  \item \textbf{Simple pole} at $t = g_*$: $\hat{Z}(t) \sim (t-g_*)^{-1}$.
        Borel sum ambiguous; imaginary part $\sim e^{-g_*/g}$ is the
        instanton contribution.
  \item \textbf{Branch point} at $t = g_*$: $\hat{Z}(t) \sim (t-g_*)^\alpha$.
        Transfer theorem gives $c_n \sim n^{-\alpha-1} n! g_*^{-n}$. The
        $n!$ growth is present but with a power-law correction. These are
        \textbf{renormalons} (IR or UV), associated with factorial growth
        from perturbation theory that is not controlled by instantons.
\end{itemize}
\end{keybox}

\subsubsection{Resurgence: The GF of All Corrections}

The modern framework of \textbf{resurgence} (\'Ecalle \cite{ecalle1981},
developed for physics by Dunne--\"Unsal \cite{dunne2014}) promotes this to
a full AC-style theory. The \textbf{trans-series}:
\[
  Z(g) \sim \sum_{k=0}^\infty e^{-k S_I/g} \sum_{n=0}^\infty c_n^{(k)} g^n
\]
where $S_I$ is the instanton action, packages the perturbative expansion
around each non-perturbative sector into a separate GF. The Borel--\'Ecalle
summation procedure is the analytic continuation of the GF past its
singularity, and the ``Stokes phenomena'' (ambiguity in the analytic
continuation) cancel between different sectors --- a non-trivial identity
that is the resurgence statement.

In AC language: the full trans-series is a GF with a \emph{multi-sheeted
Riemann surface} as its domain, and resurgence is the statement that the
GF has a consistent analytic structure across all sheets.

\begin{example}[Airy Function as a Trans-Series]
The Airy equation $y'' = zy$ has a formal divergent solution:
\[
  \mathrm{Ai}(z) \sim \frac{e^{-2z^{3/2}/3}}{2\sqrt{\pi}z^{1/4}}
  \sum_{n=0}^\infty \frac{(-1)^n \Gamma(3n+1/2)}{54^n n! \Gamma(n+1/2)} z^{-3n/2}.
\]
The $\Gamma(3n+1/2)/n!$ coefficients grow as $\sim (3n)!/(n!)^2 \sim n!\cdot C^n$,
showing factorial divergence. The AC viewpoint: $\mathrm{Ai}(z)$ is a Lagrange-type
GF evaluated at the saddle point of a cubic integral, with a square-root branch
point in the Borel plane at $t = 2/3$. The two Airy solutions $\mathrm{Ai}$
and $\mathrm{Bi}$ are the two branches of this GF, and resurgence explains
how they are related via the Stokes phenomenon.
\end{example}

\begin{resultbox}[Status]
The factorial growth / Borel singularity connection is well-known in the
resurgence literature \cite{ecalle1981, dunne2014} but phrased differently
from AC. The AC transfer theorem is not cited in the resurgence literature.
The two communities (resurgence and AC) work with the same Borel transform
but have not explicitly bridged their frameworks. This appears to be a
genuine gap in the literature.
\end{resultbox}

\subsection{The Status Map: What Is Established and What Is New}
\label{sec:statusmap}

The following table summarises the state of the art for each connection.

\begin{center}
\renewcommand{\arraystretch}{1.35}
\begin{tabular}{p{4.5cm}p{3cm}p{5cm}}
\toprule
\textbf{Connection} & \textbf{Status} & \textbf{Key References} \\
\midrule
$e^{C(z)} = Z$; symmetry factors from EGF & Classical folklore & Flajolet--Sedgewick \cite{flajolet2009}, Zinn-Justin \cite{zinjustin2002} \\
\midrule
Dyson--Schwinger = Lagrange equation & Known implicitly & Foissy \cite{foissy2007}, Yeats \cite{yeats2017} \\
Lagrange inversion generates loop expansion & Not stated explicitly & -- \\
Landau pole = Lagrange branch point & Not stated explicitly & -- \\
\midrule
Connes--Kreimer: BPHZ = Hopf antipode & Rigorously proved & Connes--Kreimer \cite{conneskreimer1998, conneskreimer2000} \\
Beta function = derivation of $H_R$ & Rigorously proved & Connes--Kreimer \cite{conneskreimer2000} \\
Renormalisation = Riemann--Hilbert problem & Rigorously proved & Connes--Kreimer \cite{conneskreimer2000} \\
\midrule
Factorial growth $\Leftrightarrow$ Borel singularity & Known (separately) & \'Ecalle \cite{ecalle1981}, Dunne--\"Unsal \cite{dunne2014} \\
Borel singularity = AC transfer theorem & Not bridged & -- \\
Trans-series = multi-sheet GF & Implicit in resurgence & \'Ecalle \cite{ecalle1981} \\
\midrule
SIR orbit = singular locus of epidemic GF & Follows from Kendall \cite{kendall1956} & Not stated explicitly \\
Doi--Peliti saddle = AC branch point & Not in literature & -- \\
SIR / Doi--Peliti / AC triangle & Not synthesised & This document \\
\bottomrule
\end{tabular}
\end{center}

\subsection{The Grand Summary: One Table}

All four connections can be summarised as a dictionary between AC and QFT:

\begin{center}
\renewcommand{\arraystretch}{1.3}
\begin{tabular}{ll}
\toprule
\textbf{Analytic Combinatorics} & \textbf{Quantum Field Theory} \\
\midrule
EGF of connected structures $\hat{C}(z)$ & Free energy $F = -\log Z$ \\
$e^{\hat{C}(z)}$ exponential formula & $Z = e^{\text{connected diagrams}}$ \\
Symmetry factor $= 1/|\mathrm{Aut}|$ & EGF overcounting $1/k!$ in SET \\
Lagrange equation $T = z\phi(T)$ & Dyson--Schwinger equation \\
Lagrange inversion coefficients & Feynman diagram enumeration \\
Branch point of Lagrange GF & Landau pole / non-pert.\ scale \\
Rooted tree Hopf algebra coproduct & BPHZ subdivergence subtraction \\
Hopf antipode & Renormalisation counterterms \\
M\"obius inversion on forests & Forest formula (BPHZ) \\
GF singularity type & Universality class / RG fixed point \\
Borel transform singularity & Instanton / renormalon \\
Transfer theorem: $n^{-\alpha-1} \cdot A^n$ & Perturbative coefficient growth \\
Multi-sheet Riemann surface & Trans-series / resurgence \\
\bottomrule
\end{tabular}
\end{center}

\begin{keybox}[Why This Matters]
The table above is not merely a curiosity. It suggests that the
\emph{combinatorial structure} of a QFT --- the grammar by which Feynman diagrams
are built --- is the primary object, and the analytic structure (propagators,
amplitudes, renormalisation) is a \emph{representation} of that combinatorial
object. This is the philosophy behind the Connes--Kreimer programme, and it
has led to concrete new results: the Birkhoff decomposition of renormalisation,
the connection to the Riemann--Hilbert problem, and new methods for computing
multi-loop amplitudes.

For the stochastic processes studied in this document (SIR, $A+A\to\emptyset$),
the same philosophy applies: the epidemic tree, the first-passage excursion,
and the Doi--Peliti path integral are different representations of the same
combinatorial grammar, and their common singularity structure is the physical
content.
\end{keybox}

\clearpage

% ========== APPENDIX ==========
\appendix

\section{Table of Key Results}
\label{app:table}

\subsection{Generating Function Identities}

\begin{center}
\begin{tabular}{lll}
\toprule
Object & GF & Notes \\
\midrule
Catalan numbers $C_n = \binom{2n}{n}/(n+1)$ & $C(z) = (1-\sqrt{1-4z})/2z$ & Binary trees \\
Borel distribution & $T(z) = ze^{R_0(T-1)}$ & SIR final size (Lagrange) \\
First-passage (1D) & $F(z) = (1-\sqrt{1-z^2})/z$ & Ballot paths \\
Poisson PGF & $g(s) = e^{\lambda(s-1)}$ & Galton--Watson offspring \\
Geometric GF & $1/(1-z)$ & Sequences of atoms \\
\bottomrule
\end{tabular}
\end{center}

\subsection{Transfer Theorem Reference}

\begin{center}
\begin{tabular}{llll}
\toprule
Singularity & Local form near $z_*$ & $[z^n]f$ behaviour & Example \\
\midrule
Simple pole & $(1-z/z_*)^{-1}$ & $\sim z_*^{-n}$ & Fibonacci \\
$k$-th pole & $(1-z/z_*)^{-k}$ & $\sim n^{k-1}z_*^{-n}$ & Path lengths \\
Sq.\ root + & $(1-z/z_*)^{1/2}$ & $\sim n^{-3/2}z_*^{-n}$ & SIR, A+A \\
Sq.\ root $-$ & $(1-z/z_*)^{-1/2}$ & $\sim n^{-1/2}z_*^{-n}$ & Simple RW \\
Log & $-\log(1-z/z_*)$ & $\sim n^{-1}z_*^{-n}$ & Cycles \\
\bottomrule
\end{tabular}
\end{center}

\subsection{Doi--Peliti Reaction Dictionary}

\begin{center}
\begin{tabular}{lll}
\toprule
Reaction & Rate & Liouvillian $\mathcal{H}$ \\
\midrule
$A \to \emptyset$ & $\delta$ & $\delta(1-\ad)a$ \\
$\emptyset \to A$ & $\sigma$ & $\sigma(\ad - 1)$ \\
$A \to 2A$ & $\lambda$ & $\lambda(\ad{}^2 - \ad)a$ \\
$2A \to \emptyset$ & $\mu$ & $\mu(1 - \ad{}^2)a^2$ \\
$2A \to A$ & $\kappa$ & $\kappa(\ad - \ad{}^2)a^2$ \\
$A + B \to \emptyset$ & $\nu$ & $\nu(1 - \ad{}_A\ad{}_B)a_Ab_B$ \\
\bottomrule
\end{tabular}
\end{center}

The general rule: reaction $kA \to \ell A$ at rate $r$ gives
$\mathcal{H} = r(\ad{}^\ell - \ad{}^k)a^k$.

\subsection{Critical Dimensions for Reaction--Diffusion}

\begin{center}
\begin{tabular}{llll}
\toprule
Reaction & Upper $d_c$ & Exponent $\alpha$ ($d<d_c$) & Mean-field ($d>d_c$) \\
\midrule
$A + A \to \emptyset$ & 2 & $d/2$ & 1 \\
$A + A \to A$ & 2 & $d/2$ & 1 \\
$A + B \to \emptyset$ & 4 & $d/4$ & 1 \\
$3A \to \emptyset$ & 1 & $d/2$ & 1/2 \\
\bottomrule
\end{tabular}
\end{center}

\subsection{The Weyl Algebra and the GF Isomorphism}

The \textbf{Weyl algebra} $W_1$ is the unital associative algebra generated by $p$
and $q$ with relation $[p, q] = 1$ (or equivalently, $pq - qp = 1$). Under the
representation $p = \partial_z$, $q = z$ (multiplication), $W_1$ acts on $\C[[z]]$.
This is the unique irreducible representation of $W_1$ on a space of formal
power series (Stone--von Neumann theorem in the algebraic sense).

Doi's identification: $a \equiv \partial_z$, $\ad \equiv z$ gives an isomorphism:
\begin{itemize}[leftmargin=2em]
  \item Fock space $\mathcal{F} \cong \C[[z]]$ as $W_1$-modules.
  \item The state $|\psi\rangle = \sum_n P_n |n\rangle$ corresponds to $G(z) = \sum_n P_n z^n$.
  \item The Liouvillian $\mathcal{H}(a, \ad)$ corresponds to $\mathcal{H}(\partial_z, z)$.
  \item The evolution equation $\partial_t|\psi\rangle = \mathcal{H}|\psi\rangle$
        becomes the GF PDE $\partial_t G = \mathcal{H}(\partial_z, z)G$.
\end{itemize}

\subsection{The Hopf Algebra Connection}

The coproduct of $W_1$ (describing how two particles interact) is:

\[
  \Delta(a) = a \otimes 1 + 1 \otimes a, \qquad
  \Delta(\ad) = \ad \otimes 1 + 1 \otimes \ad.
\]

The combinatorial counterpart: the coproduct of the Connes--Kreimer Hopf algebra
of rooted trees decomposes a tree into subtree/quotient pairs, which is precisely
the admissible cut operation. Renormalisation (the antipode of the Hopf algebra)
corresponds to the inclusion--exclusion in Flajolet's symbolic method for the
MSET and CYC constructions. This is why the Feynman diagram expansion of the
Doi--Peliti theory and the AC generating function equations are algebraically
the same object at different levels of abstraction.

\section{Literature Guide}
\label{app:lit}

\subsection{Stochastic SIR}

\begin{itemize}[leftmargin=2em]
  \item \textbf{Kermack \& McKendrick (1927)} \cite{kermack1927}: the original
        deterministic SIR paper. Derives the epidemic threshold and conservation law.

  \item \textbf{Whittle (1955)} \cite{whittle1955}: first use of the PGF for
        the stochastic SIR. Derives the minor-outbreak probability.

  \item \textbf{Kendall (1956)} \cite{kendall1956}: explicit identification of the
        characteristic equations of the PGF PDE with the deterministic SIR.

  \item \textbf{Chalub \& Souza (2011)} \cite{chalub2011}: modern PDE formulation.
        Proves that the SIR ODEs are the characteristics and derives large-population
        limits rigorously.
\end{itemize}

\subsection{Doi--Peliti}

\begin{itemize}[leftmargin=2em]
  \item \textbf{Doi (1976a,b)} \cite{doi1976a,doi1976b}: two papers introducing the
        second-quantisation formalism for reaction--diffusion master equations.

  \item \textbf{Peliti (1985)} \cite{peliti1985}: coherent-state path integral
        for the Doi formalism. Introduces the $\phi, \tilde\phi$ field doubling.

  \item \textbf{Lee (1994)} \cite{lee1994}: RG calculation for $kA\to\emptyset$.
        Proves that the density exponent is exact.

  \item \textbf{T\"auber, Howard \& Vollmayr-Lee (2005)} \cite{tauber2005}: the
        comprehensive review. Covers $kA\to\ell A$, $A+B\to\emptyset$, disorder,
        and absorbing-state transitions. Best entry point.

  \item \textbf{T\"auber (2014)} \cite{tauber2014}: textbook treatment, Chapter 9
        on reaction--diffusion.
\end{itemize}

\subsection{Analytic Combinatorics}

\begin{itemize}[leftmargin=2em]
  \item \textbf{Flajolet \& Sedgewick (2009)} \cite{flajolet2009}: the canonical
        reference. Parts A--B cover the symbolic method; Part C covers singularity
        analysis and transfer theorems. Freely available online.

  \item \textbf{Wilf (2006)} \cite{wilf2006}: \emph{generatingfunctionology}, a
        more accessible introduction to generating functions in combinatorics.

  \item \textbf{Pemantle \& Wilson (2013)} \cite{pemantle2013}: multivariate
        analytic combinatorics; extends the singularity approach to GFs in multiple
        variables.
\end{itemize}

\subsection{AC, Feynman Diagrams, and Renormalisation}

\begin{itemize}[leftmargin=2em]
  \item \textbf{Connes \& Kreimer (1998)} \cite{conneskreimer1998}: the discovery
        that the BPHZ forest formula is the antipode of a Hopf algebra of rooted
        trees. The paper that started the programme. Dense but repays careful reading.

  \item \textbf{Connes \& Kreimer (2000)} \cite{conneskreimer2000}: the
        Riemann--Hilbert problem formulation of renormalisation. Shows the beta
        function is a derivation of the Hopf algebra and that the Birkhoff
        decomposition encodes counterterms.

  \item \textbf{Yeats (2017)} \cite{yeats2017}: \emph{A Combinatorial Perspective
        on Quantum Field Theory}. The best entry point for a mathematician.
        Explicitly begins from perturbative expansions as generating functions,
        introduces the Connes--Kreimer algebra, and works through Dyson--Schwinger
        equations as combinatorial fixed-point equations. The table of contents
        mirrors the structure of Section~\ref{sec:feynman} of this document.

  \item \textbf{Foissy (2007)} \cite{foissy2007}: combinatorial Dyson--Schwinger
        equations as systems of equations in the Hopf algebra of planar trees.
        This is the closest existing work to the ``Lagrange equation = DSE''
        identification in Section~\ref{sec:dse_lagrange}.

  \item \textbf{Manchon (2004)} \cite{manchon2004}: a concise mathematical survey
        of Hopf algebras in renormalisation, including the connection to
        M\"obius inversion and the BPHZ formula.

  \item \textbf{Kreimer (2000)}: \emph{Combinatorics of (perturbative) Quantum
        Field Theory}, arXiv:hep-th/0010059. A readable overview of how Feynman
        graphs map to rooted trees and how the Hopf algebra coproduct encodes
        subdivergences.
\end{itemize}

\subsection{Resurgence and Borel Summation}

\begin{itemize}[leftmargin=2em]
  \item \textbf{\'Ecalle (1981)} \cite{ecalle1981}: the original mathematical
        framework of resurgence and analysable functions. Very technical.
        The three-volume work that introduced trans-series.

  \item \textbf{Dunne \& \"Unsal (2014)} \cite{dunne2014}: the clearest modern
        physics treatment of resurgence with explicit QFT examples. Introduces
        the connection between IR renormalons (Borel singularities) and
        non-perturbative contributions.

  \item \textbf{Zinn-Justin (2002)} \cite{zinjustin2002}: standard QFT textbook
        with good treatment of the large-order behaviour of perturbation theory,
        instantons, and Borel summation. Chapter 37--40 are directly relevant.
\end{itemize}

\clearpage
\bibliographystyle{unsrt}
\begin{thebibliography}{99}

\bibitem{kermack1927}
W.\ O.\ Kermack and A.\ G.\ McKendrick,
\newblock A contribution to the mathematical theory of epidemics,
\newblock \emph{Proc.\ R.\ Soc.\ London A}, 115:700--721, 1927.

\bibitem{whittle1955}
P.\ Whittle,
\newblock The outcome of a stochastic epidemic---a note on Bailey's paper,
\newblock \emph{Biometrika}, 42:116--122, 1955.

\bibitem{kendall1956}
D.\ G.\ Kendall,
\newblock Deterministic and stochastic epidemics in closed populations,
\newblock \emph{Proc.\ Third Berkeley Symp.\ Math.\ Statist.\ Probab.},
  4:149--165, 1956.

\bibitem{chalub2011}
F.\ A.\ C.\ C.\ Chalub and M.\ O.\ Souza,
\newblock The {SIR} epidemic model from a {PDE} point of view,
\newblock \emph{Math.\ Comput.\ Modelling}, 53(7--8):1568--1574, 2011.

\bibitem{doi1976a}
M.\ Doi,
\newblock Stochastic theory of diffusion-controlled reaction,
\newblock \emph{J.\ Phys.\ A: Math.\ Gen.}, 9:1465--1477, 1976.

\bibitem{doi1976b}
M.\ Doi,
\newblock Second quantization representation for classical many-particle system,
\newblock \emph{J.\ Phys.\ A: Math.\ Gen.}, 9:1479--1495, 1976.

\bibitem{peliti1985}
L.\ Peliti,
\newblock Path integral approach to birth-death processes on a lattice,
\newblock \emph{J.\ Phys.\ (Paris)}, 46:1469--1483, 1985.

\bibitem{lee1994}
B.\ P.\ Lee,
\newblock Renormalization group calculation for the reaction $kA\to\emptyset$,
\newblock \emph{J.\ Phys.\ A: Math.\ Gen.}, 27:2633--2652, 1994.

\bibitem{tauber2005}
U.\ C.\ T\"auber, M.\ Howard, and B.\ P.\ Vollmayr-Lee,
\newblock Applications of field-theoretic renormalization group methods to
  reaction--diffusion problems,
\newblock \emph{J.\ Phys.\ A: Math.\ Gen.}, 38:R79--R185, 2005.
\newblock arXiv:cond-mat/0501678.

\bibitem{tauber2014}
U.\ C.\ T\"auber,
\newblock \emph{Critical Dynamics},
\newblock Cambridge University Press, 2014.
\newblock Chapter 9: Reaction--Diffusion Systems.

\bibitem{flajolet2009}
P.\ Flajolet and R.\ Sedgewick,
\newblock \emph{Analytic Combinatorics},
\newblock Cambridge University Press, 2009.
\newblock Available at \url{http://algo.inria.fr/flajolet/Publications/book.pdf}.

\bibitem{wilf2006}
H.\ S.\ Wilf,
\newblock \emph{generatingfunctionology},
\newblock 3rd ed., A.\ K.\ Peters, 2006.
\newblock Available at \url{https://www.math.upenn.edu/~wilf/gfologyLinked2.pdf}.

\bibitem{pemantle2013}
R.\ Pemantle and M.\ C.\ Wilson,
\newblock \emph{Analytic Combinatorics in Several Variables},
\newblock Cambridge University Press, 2013.

\bibitem{conneskreimer1998}
A.\ Connes and D.\ Kreimer,
\newblock Hopf algebras, renormalization and noncommutative geometry,
\newblock \emph{Commun.\ Math.\ Phys.}, 199:203--242, 1998.
\newblock arXiv:hep-th/9808042.

\bibitem{conneskreimer2000}
A.\ Connes and D.\ Kreimer,
\newblock Renormalization in quantum field theory and the {Riemann}--{Hilbert}
  problem {I}: the {Hopf} algebra structure of graphs and the main theorem,
\newblock \emph{Commun.\ Math.\ Phys.}, 210:249--273, 2000.
\newblock arXiv:hep-th/9912092.

\bibitem{yeats2017}
K.\ Yeats,
\newblock \emph{A Combinatorial Perspective on Quantum Field Theory},
\newblock SpringerBriefs in Mathematical Physics 15, Springer, 2017.

\bibitem{foissy2007}
L.\ Foissy,
\newblock Fa\`a di Bruno subalgebras of the {Hopf} algebra of planar trees
  from combinatorial {Dyson}--{Schwinger} equations,
\newblock \emph{Adv.\ Math.}, 218(1):136--162, 2008.
\newblock arXiv:0707.1204.

\bibitem{manchon2004}
D.\ Manchon,
\newblock Hopf algebras in renormalisation,
\newblock In \emph{Handbook of Algebra}, vol.\ 5, pp.\ 365--427,
  Elsevier, 2008.
\newblock arXiv:math/0408405.

\bibitem{ecalle1981}
J.\ \'Ecalle,
\newblock \emph{Les Fonctions R\'esurgentes}, Vols.\ 1--3,
\newblock Publ.\ Math.\ d'Orsay, 1981.

\bibitem{dunne2014}
G.\ V.\ Dunne and M.\ \"Unsal,
\newblock Uniform {WKB}, multi-instantons, and resurgent trans-series,
\newblock \emph{Phys.\ Rev.\ D}, 89:105009, 2014.
\newblock arXiv:1401.5202.

\bibitem{zinjustin2002}
J.\ Zinn-Justin,
\newblock \emph{Quantum Field Theory and Critical Phenomena},
\newblock 4th ed., Oxford University Press, 2002.

\end{thebibliography}

\end{document}
